% ══════════════════════════════════════════════════════════════════════════════
%                         فصل ششم
%                 آنارشیسم و چپ آزادی‌خواه
% ══════════════════════════════════════════════════════════════════════════════

\chapter{آنارشیسم و چپ آزادی‌خواه}
\label{ch:anarchism}

\begin{flushright}
\textit{«آنارشی مادر نظم است.»}

— پیر ژوزف پرودون
\end{flushright}

% ══════════════════════════════════════════════════════════════════════════════
%                         مقدمه فصل
% ══════════════════════════════════════════════════════════════════════════════

\lettrine[lines=3, lraise=0.1, nindent=0.5em]{\textcolor{OrangeAccent}{آ}}{نارشیسم} 
شاید بدفهم‌ترین ایدئولوژی سیاسی باشد. در ذهن عموم، آنارشیسم با هرج‌ومرج، خشونت و بی‌نظمی تداعی می‌شود. اما آنارشیسم واقعی، نه نفی نظم، بلکه نفی \textbf{سلطه} است. آنارشیست‌ها می‌گویند: انسان‌ها می‌توانند بدون دولت، بدون سلسله‌مراتب اجباری، و بدون استثمار زندگی کنند—و این نه هرج‌ومرج، بلکه بالاترین شکل نظم است.

% ══════════════════════════════════════════════════════════════════════════════
\section{آنارشیسم چیست؟}
% ══════════════════════════════════════════════════════════════════════════════

\subsection{تعریف و اصول بنیادین}

\begin{defbox}[تعریف: آنارشیسم]
از واژه یونانی \textit{an-archos} به معنای «بدون حاکم» یا «بدون فرمانروا».

\textbf{اصول مشترک اغلب آنارشیست‌ها:}
\begin{itemize}[nosep]
    \item \textbf{ضد دولت:} دولت ذاتاً سرکوبگر است، حتی دولت «کارگری»
    \item \textbf{ضد سلسله‌مراتب:} هر شکل از سلطه اجباری باید به چالش کشیده شود
    \item \textbf{خودگردانی:} انسان‌ها می‌توانند خودشان امورشان را اداره کنند
    \item \textbf{کمک متقابل:} همکاری داوطلبانه جای رقابت و اجبار را می‌گیرد
    \item \textbf{اقدام مستقیم:} تغییر از پایین، نه از طریق پارلمان یا حزب
\end{itemize}
\end{defbox}

\begin{debatebox}[آنارشیسم ≠ هرج‌ومرج]
\textbf{تصور غلط:} آنارشیست‌ها خواهان هرج‌ومرج و بی‌قانونی‌اند.

\textbf{پاسخ آنارشیست‌ها:}
\begin{itemize}[nosep]
    \item ما خواهان \textbf{نظم بدون سلطه} هستیم
    \item قوانین می‌توانند توسط جامعه تعیین شوند، نه دولت
    \item بسیاری از «بی‌نظمی‌ها» محصول خودِ دولت و سرمایه‌داری است
    \item جوامع بی‌دولت در تاریخ وجود داشته‌اند
\end{itemize}

\textbf{نقد منتقدان:}
\begin{itemize}[nosep]
    \item چگونه بدون دولت از حقوق دفاع می‌شود؟
    \item چه کسی جلوی زورگویی را می‌گیرد؟
    \item آیا این آرمان‌شهرگرایی محض نیست؟
\end{itemize}
\end{debatebox}

\subsection{تفاوت با سایر جریان‌های چپ}

\begin{table}[htbp]
\centering
\caption{مقایسه آنارشیسم با مارکسیسم و لیبرالیسم}
\label{tab:anarchism-compare}
\begin{tabular}{r|c|c|c}
\toprule
\textbf{موضوع} & \textbf{آنارشیسم} & \textbf{مارکسیسم} & \textbf{لیبرالیسم} \\
\midrule
دولت & باید فوراً لغو شود & ابزار گذار، سپس زوال & لازم اما محدود \\
انقلاب & از پایین، بی‌حزب & رهبری حزب پیشاهنگ & اصلاحات تدریجی \\
مالکیت & اشتراکی یا فردیِ کار & دولتی/اشتراکی & خصوصی \\
آزادی & هدف اصلی و فوری & پس از انقلاب & در چارچوب قانون \\
سازمان‌دهی & افقی، فدرالی & متمرکز، سلسله‌مراتبی & نمایندگی \\
\bottomrule
\end{tabular}
\end{table}

% ══════════════════════════════════════════════════════════════════════════════
\section{پیشگامان آنارشیسم}
% ══════════════════════════════════════════════════════════════════════════════

\subsection{ویلیام گادوین}

\begin{personbox}[ویلیام گادوین (۱۷۵۶-۱۸۳۶)]

\textbf{ملیت:} انگلیسی

\textbf{اثر کلیدی:} \book{تحقیق در باب عدالت سیاسی} (۱۷۹۳)

\textbf{ایده محوری:} عقل‌گرایی آنارشیستی

\mediumspace

\person{گادوین} اغلب «پدر آنارشیسم فلسفی» نامیده می‌شود:

\begin{itemize}[nosep]
    \item دولت، قانون و حتی ازدواج نهادهای سرکوبگرند
    \item با گسترش \textbf{عقل و آموزش}، نیاز به دولت از بین می‌رود
    \item مالکیت خصوصی مانع عدالت است
    \item روش: اقناع عقلانی، نه خشونت
\end{itemize}

\mediumspace

\textbf{همسر:} مری وولستون‌کرافت (فمینیست پیشگام)

\textbf{دختر:} مری شلی (نویسنده فرانکنشتاین)
\end{personbox}

\subsection{پیر ژوزف پرودون}

\begin{personbox}[پیر ژوزف پرودون (۱۸۰۹-۱۸۶۵)]

\textbf{ملیت:} فرانسوی

\textbf{آثار کلیدی:} \book{مالکیت چیست؟} (۱۸۴۰)، \book{فلسفه فقر}

\textbf{جمله معروف:} «مالکیت دزدی است!» و «آنارشی مادر نظم است»

\mediumspace

\person{پرودون} نخستین کسی بود که خود را \textbf{آنارشیست} نامید:

\begin{itemize}[nosep]
    \item نقد مالکیت بزرگ (اجاره، سود، بهره = دزدی)
    \item اما دفاع از «تصرف» و مالکیت شخصی کوچک
    \item \textbf{موتوئالیسم:} تعاونی‌ها و بانک خلق (اعتبار بدون بهره)
    \item فدرالیسم: جوامع کوچک خودگردان
    \item مخالفت با انقلاب خشونت‌آمیز
\end{itemize}

\mediumspace

\textbf{نقد مارکس:} مارکس در \book{فقر فلسفه} پرودون را نقد کرد و او را خرده‌بورژوا خواند.
\end{personbox}

\subsection{ماکس اشتیرنر}

\begin{personbox}[ماکس اشتیرنر (۱۸۰۶-۱۸۵۶)]

\textbf{نام اصلی:} یوهان کاسپار اشمیت

\textbf{ملیت:} آلمانی

\textbf{اثر کلیدی:} \book{یگانه و مالِ او} (۱۸۴۴)

\textbf{ایده محوری:} فردگرایی رادیکال

\mediumspace

\person{اشتیرنر} آنارشیست فردگرای افراطی بود:

\begin{itemize}[nosep]
    \item همه مفاهیم انتزاعی (دولت، خدا، انسانیت، جامعه) «شبح»اند
    \item فرد تنها واقعیت است
    \item «من» ارزش‌های خودم را می‌سازم
    \item هیچ قانون یا اخلاقی بر من تحمیل نیست
    \item «اتحادیه خودخواهان» به جای جامعه
\end{itemize}

\mediumspace

\textbf{تأثیر:} پیش‌درآمد نیچه، اگزیستانسیالیسم، برخی لیبرتارین‌های راست
\end{personbox}

% ══════════════════════════════════════════════════════════════════════════════
\section{آنارشیسم جمعی: باکونین}
% ══════════════════════════════════════════════════════════════════════════════

\begin{personbox}[میخائیل باکونین (۱۸۱۴-۱۸۷۶)]

\textbf{ملیت:} روسی

\textbf{آثار کلیدی:} \book{دولت و آنارشی}، \book{خدا و دولت}

\textbf{ایده محوری:} آنارشیسم جمعی و انقلابی

\mediumspace

\person{باکونین} آنارشیسم را به جنبش کارگری پیوند زد و رقیب اصلی مارکس در انترناسیونال اول شد:

\begin{itemize}[nosep]
    \item انقلاب باید \textbf{همزمان} دولت و سرمایه‌داری را نابود کند
    \item نباید دولت جدید (حتی کارگری) ساخت
    \item \textbf{آزادی جمعی:} آزادی من بدون آزادی همه کامل نیست
    \item سازمان‌دهی فدرالی و از پایین
    \item لومپن‌پرولتاریا (حاشیه‌نشینان) هم انقلابی‌اند
\end{itemize}

\mediumspace

\textbf{نقل‌قول معروف:} \textit{«اگر خدا وجود داشت، باید او را نابود کرد.»}

\textbf{پیش‌بینی درباره مارکسیسم:} \textit{«اگر طبقه کارگر قدرت را بگیرد، به بوروکراسی جدیدی تبدیل می‌شود که بر توده‌ها حکومت می‌کند.»}
\end{personbox}

\begin{figure}[htbp]
\centering
\begin{tikzpicture}[scale=0.85]

% عنوان
\node[font=\large\bfseries, text=DarkGray] at (0, 6.5) {\fa{جدال مارکس و باکونین در انترناسیونال اول}};

% کادر مارکس
\fill[LeftRed, opacity=0.2, rounded corners=10pt] (-7, 0) rectangle (-0.5, 5.5);
\draw[line width=2pt, color=LeftRed, rounded corners=10pt] (-7, 0) rectangle (-0.5, 5.5);

\node[text=LeftRed, font=\bfseries\large] at (-3.75, 5) {\fa{مارکس}};

\node[text=DarkGray, font=\small, text width=5.5cm, align=center] at (-3.75, 3) {
    \fa{دولت کارگری لازم است}\\[5pt]
    \fa{ابتدا قدرت، سپس تغییر}\\[5pt]
    \fa{حزب متمرکز}\\[5pt]
    \fa{پرولتاریای صنعتی}\\[5pt]
    \fa{علم و نظریه مهم است}
};

\node[text=LeftRedDark, font=\small\bfseries] at (-3.75, 0.7) {\fa{دیکتاتوری پرولتاریا}};

% کادر باکونین
\fill[OrangeAccent, opacity=0.2, rounded corners=10pt] (0.5, 0) rectangle (7, 5.5);
\draw[line width=2pt, color=OrangeAccent, rounded corners=10pt] (0.5, 0) rectangle (7, 5.5);

\node[text=OrangeAccent, font=\bfseries\large] at (3.75, 5) {\fa{باکونین}};

\node[text=DarkGray, font=\small, text width=5.5cm, align=center] at (3.75, 3) {
    \fa{هر دولتی سرکوبگر است}\\[5pt]
    \fa{نابودی همزمان دولت و سرمایه}\\[5pt]
    \fa{فدراسیون از پایین}\\[5pt]
    \fa{دهقانان و لومپن‌ها هم مهم‌اند}\\[5pt]
    \fa{عمل مستقیم}
};

\node[text=OrangeAccent!80!black, font=\small\bfseries] at (3.75, 0.7) {\fa{فدراسیون آزاد}};

% علامت تقابل
\node[text=GoldAccent, font=\Huge\bfseries] at (0, 2.75) {\fa{⚔}};

% نتیجه
\node[text=MediumGray, font=\small, text width=12cm, align=center] at (0, -0.8) {
    \fa{نتیجه: اخراج باکونین از انترناسیونال (۱۸۷۲) و انشعاب جنبش سوسیالیستی}
};

\end{tikzpicture}
\caption{مقایسه مواضع مارکس و باکونین}
\label{fig:marx-bakunin}
\end{figure}

% ══════════════════════════════════════════════════════════════════════════════
\section{آنارکو-کمونیسم: کروپوتکین}
% ══════════════════════════════════════════════════════════════════════════════

\begin{personbox}[پیوتر کروپوتکین (۱۸۴۲-۱۹۲۱)]

\textbf{ملیت:} روسی (شاهزاده!)

\textbf{آثار کلیدی:} \book{کمک متقابل} (۱۹۰۲)، \book{تسخیر نان} (۱۸۹۲)

\textbf{ایده محوری:} آنارکو-کمونیسم و کمک متقابل

\mediumspace

\person{کروپوتکین}، زیست‌شناس و جغرافی‌دان، آنارشیسم را با علم پیوند زد:

\begin{itemize}[nosep]
    \item \textbf{کمک متقابل:} همکاری، نه رقابت، موتور تکامل است
    \item داروینیسم اجتماعی (بقای اصلح) علم نیست، ایدئولوژی است
    \item \textbf{آنارکو-کمونیسم:} مالکیت اشتراکی + توزیع بر اساس نیاز
    \item نه دستمزد، نه پول، نه بازار
    \item «از هرکس به توانش، به هرکس به نیازش»—فوراً، نه پس از گذار
\end{itemize}

\mediumspace

\textbf{تفاوت با باکونین:} باکونین «جمع‌گرا» بود (توزیع به نسبت کار)؛ کروپوتکین «کمونیست» (توزیع به نسبت نیاز).

\mediumspace

\textbf{نقل‌قول:} \textit{«ما ثروتمندیم، بسیار ثروتمند. اما چرا بسیاری فقیرند؟»}
\end{personbox}

\begin{defbox}[تعریف: کمک متقابل]
نظریه کروپوتکین که بر اساس مشاهدات علمی استدلال می‌کند:
\begin{itemize}[nosep]
    \item در طبیعت، گونه‌هایی که \textbf{همکاری} می‌کنند موفق‌ترند
    \item جوامع انسانی تاریخ طولانی کمک متقابل دارند (روستا، صنف، تعاونی)
    \item سرمایه‌داری و دولت این غریزه طبیعی را سرکوب کرده‌اند
    \item آنارشی بازگشت به همکاری طبیعی است، نه آرمان‌شهر
\end{itemize}
\end{defbox}

% ══════════════════════════════════════════════════════════════════════════════
\section{آنارکو-سندیکالیسم}
% ══════════════════════════════════════════════════════════════════════════════

آنارکو-سندیکالیسم، آنارشیسم را با جنبش اتحادیه‌ای ترکیب کرد و در اوایل قرن بیستم جنبشی توده‌ای شد.

\begin{defbox}[تعریف: آنارکو-سندیکالیسم]
\begin{itemize}[nosep]
    \item \textbf{سندیکا (اتحادیه کارگری)} ابزار اصلی مبارزه و سازمان جامعه آینده
    \item اقدام مستقیم: اعتصاب، بایکوت، خرابکاری
    \item \textbf{اعتصاب عمومی} می‌تواند سرمایه‌داری را فلج کند
    \item پس از انقلاب: اتحادیه‌ها کارخانه‌ها را اداره می‌کنند
    \item نه حزب سیاسی، نه دولت
\end{itemize}
\end{defbox}

\begin{historybox}[جنبش‌های آنارکو-سندیکالیست]
\begin{itemize}[nosep]
    \item \textbf{CGT فرانسه:} اتحادیه عمومی کار، تأثیرگذارترین در اوایل قرن
    \item \textbf{IWW آمریکا:} «کارگران صنعتی جهان» (Wobblies)، سازمان‌دهی همه کارگران
    \item \textbf{CNT اسپانیا:} کنفدراسیون ملی کار، میلیون‌ها عضو
    \item \textbf{FORA آرژانتین:} فدراسیون کارگران منطقه‌ای
\end{itemize}
\end{historybox}

\subsection{انقلاب اسپانیا (۱۹۳۶-۱۹۳۹)}

\begin{figure}[htbp]
\centering
\begin{tikzpicture}[scale=0.85]

% عنوان
\node[font=\large\bfseries, text=DarkGray] at (0, 7) {\fa{تایم‌لاین: انقلاب اسپانیا}};

% خط زمان
\draw[line width=4pt, MediumGray] (-7, 5) -- (7, 5);

% رویدادها
% انتخابات
\fill[CenterGreen] (-6, 5) circle (9pt);
\node[above, text=CenterGreen, font=\scriptsize\bfseries] at (-6, 5.4) {\fa{فوریه ۳۶}};
\node[below, text width=2cm, align=center, text=DarkGray, font=\tiny] at (-6, 4.4) {
    \fa{پیروزی}\\
    \fa{جبهه مردمی}
};

% کودتا
\fill[RightBlueDark] (-3.5, 5) circle (9pt);
\node[above, text=RightBlueDark, font=\scriptsize\bfseries] at (-3.5, 5.4) {\fa{ژوئیه ۳۶}};
\node[below, text width=2cm, align=center, text=DarkGray, font=\tiny] at (-3.5, 4.4) {
    \fa{کودتای فرانکو}\\
    \fa{جنگ داخلی}
};

% انقلاب
\fill[OrangeAccent] (-1, 5) circle (10pt);
\node[above, text=OrangeAccent, font=\scriptsize\bfseries] at (-1, 5.4) {\fa{ژوئیه-اوت ۳۶}};
\node[below, text width=2.3cm, align=center, text=DarkGray, font=\tiny] at (-1, 4.3) {
    \fa{انقلاب آنارشیستی}\\
    \fa{در کاتالونیا}\\
    \fa{و آراگون}
};

% سرکوب
\fill[LeftRed] (2, 5) circle (9pt);
\node[above, text=LeftRed, font=\scriptsize\bfseries] at (2, 5.4) {\fa{مه ۳۷}};
\node[below, text width=2cm, align=center, text=DarkGray, font=\tiny] at (2, 4.4) {
    \fa{درگیری بارسلونا}\\
    \fa{سرکوب توسط کمونیست‌ها}
};

% شکست
\fill[MediumGray] (5, 5) circle (9pt);
\node[above, text=MediumGray, font=\scriptsize\bfseries] at (5, 5.4) {\fa{آوریل ۳۹}};
\node[below, text width=2cm, align=center, text=DarkGray, font=\tiny] at (5, 4.4) {
    \fa{پیروزی فرانکو}\\
    \fa{پایان جمهوری}
};

% توضیح
\node[text=OrangeAccent, font=\small\bfseries, text width=12cm, align=center] at (0, 2.8) {
    \fa{در تابستان ۱۹۳۶، کارگران و دهقانان کاتالونیا و آراگون}\\
    \fa{کارخانه‌ها و زمین‌ها را تصرف کردند—بزرگ‌ترین تجربه آنارشیسم در تاریخ.}
};

\end{tikzpicture}
\caption{سیر رویدادهای انقلاب اسپانیا}
\label{fig:spanish-revolution}
\end{figure}

\begin{historybox}[انقلاب اجتماعی در اسپانیا]
در مناطق تحت کنترل جمهوری‌خواهان، به‌ویژه کاتالونیا:
\begin{itemize}[nosep]
    \item کارخانه‌ها توسط \textbf{کمیته‌های کارگری} اداره شدند
    \item زمین‌ها \textbf{اشتراکی} شدند (در برخی روستاها پول لغو شد)
    \item زنان نقش فعال گرفتند (\textbf{موخرس لیبرس}—زنان آزاد)
    \item میلیشیای مردمی جای ارتش را گرفت
    \item جورج اورول در کتاب \book{ادای احترام به کاتالونیا} این تجربه را توصیف کرد
\end{itemize}

\textbf{سرنوشت:} فشار کمونیست‌ها (با حمایت شوروی) برای «نظم» و متمرکزسازی، و سرانجام پیروزی فاشیست‌ها.
\end{historybox}

\begin{quotebox}[جورج اورول، ادای احترام به کاتالونیا]
\textit{«برای نخستین بار در شهری بودم که طبقه کارگر در مسند قدرت بود. تقریباً هر ساختمانی توسط کارگران تصرف شده بود... هر دیوار با داس‌وچکش و نام‌های حزب انقلابی پوشیده بود... گارسون‌ها و فروشندگان مستقیم به چشمانت نگاه می‌کردند و با تو برابر رفتار می‌کردند.»}
\end{quotebox}

% ══════════════════════════════════════════════════════════════════════════════
\section{آنارشیسم فردگرا}
% ══════════════════════════════════════════════════════════════════════════════

در کنار آنارشیسم اجتماعی، گرایش فردگرا نیز وجود داشته است.

\begin{personbox}[بنجامین تاکر (۱۸۵۴-۱۹۳۹)]

\textbf{ملیت:} آمریکایی

\textbf{نشریه:} \textit{لیبرتی} (۱۸۸۱-۱۹۰۸)

\textbf{ایده محوری:} آنارشیسم فردگرای بازار

\mediumspace

\person{تاکر} گرایش آمریکایی آنارشیسم را نمایندگی می‌کرد:

\begin{itemize}[nosep]
    \item تأثیرپذیری از پرودون و اشتیرنر
    \item \textbf{چهار انحصار} باید لغو شوند: پول، زمین، تعرفه، پتنت
    \item با لغو این انحصارها، بازار آزاد واقعی شکل می‌گیرد
    \item کارگران می‌توانند سرمایه‌داران را شکست دهند
    \item مخالف کمونیسم و اشتراک اجباری
\end{itemize}

\mediumspace

\textbf{میراث:} برخی او را پیش‌درآمد آنارکو-کاپیتالیسم می‌دانند، برخی نه.
\end{personbox}

\begin{debatebox}[آنارکو-کاپیتالیسم: آنارشیسم یا نه؟]
در دهه ۱۹۶۰، \person{موری راتبارد} «آنارکو-کاپیتالیسم» را مطرح کرد:
\begin{itemize}[nosep]
    \item لغو دولت + حفظ سرمایه‌داری و مالکیت خصوصی
    \item امنیت و دادگستری هم خصوصی شوند
    \item بازار آزاد راه‌حل همه مشکلات است
\end{itemize}

\textbf{آنارشیست‌های سنتی:}
\begin{itemize}[nosep]
    \item این آنارشیسم نیست—آنارشیسم ضد سرمایه‌داری است
    \item سرمایه‌داری بدون دولت یعنی سلطه شرکت‌ها
    \item نام «لیبرتارین» را از ما دزدیده‌اند
\end{itemize}

\textbf{نتیجه:} اختلاف عمیق و مصالحه‌ناپذیر.
\end{debatebox}

% ══════════════════════════════════════════════════════════════════════════════
\section{جریان‌های معاصر آنارشیسم}
% ══════════════════════════════════════════════════════════════════════════════

آنارشیسم در دهه‌های اخیر احیا شده و شاخه‌های جدیدی روییده است.

\begin{figure}[htbp]
\centering
\begin{tikzpicture}[
    grow=down,
    level 1/.style={sibling distance=4cm, level distance=2.2cm},
    level 2/.style={sibling distance=2.5cm, level distance=1.8cm},
    edge from parent/.style={draw, line width=1.5pt, color=OrangeAccent},
    every node/.style={font=\small}
]

% ریشه
\node[rectangle, rounded corners=5pt, fill=OrangeAccent, text=white, font=\bfseries, minimum width=3.5cm, minimum height=0.8cm] {\fa{آنارشیسم معاصر}}
    child {
        node[rectangle, rounded corners=5pt, fill=OrangeAccent!70, text=white, minimum width=2.5cm, minimum height=0.6cm] {\fa{اجتماعی}}
        child {
            node[rectangle, rounded corners=3pt, fill=LeftRedLight, text=white, minimum width=2cm, font=\scriptsize] {\fa{آنارکو-کمونیسم}}
        }
        child {
            node[rectangle, rounded corners=3pt, fill=LeftRed, text=white, minimum width=2cm, font=\scriptsize] {\fa{سندیکالیسم}}
        }
    }
    child {
        node[rectangle, rounded corners=5pt, fill=CenterGreen, text=white, minimum width=2.5cm, minimum height=0.6cm] {\fa{سبز}}
        child {
            node[rectangle, rounded corners=3pt, fill=CenterGreenLight, text=DarkGray, minimum width=2cm, font=\scriptsize] {\fa{اکو-آنارشیسم}}
        }
        child {
            node[rectangle, rounded corners=3pt, fill=CenterGreenDark, text=white, minimum width=2cm, font=\scriptsize] {\fa{بدوی‌گرایی}}
        }
    }
    child {
        node[rectangle, rounded corners=5pt, fill=PurpleAccent, text=white, minimum width=2.5cm, minimum height=0.6cm] {\fa{هویتی}}
        child {
            node[rectangle, rounded corners=3pt, fill=PurpleLight, text=DarkGray, minimum width=2cm, font=\scriptsize] {\fa{آنارکا-فمینیسم}}
        }
        child {
            node[rectangle, rounded corners=3pt, fill=PurpleAccent!70, text=white, minimum width=2cm, font=\scriptsize] {\fa{کوئیر}}
        }
    };

% عنوان
\node[font=\large\bfseries, text=DarkGray] at (0, 2) {\fa{شاخه‌های آنارشیسم معاصر}};

\end{tikzpicture}
\caption{انواع آنارشیسم معاصر}
\label{fig:anarchism-branches}
\end{figure}

\subsection{آنارکا-فمینیسم}

\begin{defbox}[تعریف: آنارکا-فمینیسم]
ترکیب آنارشیسم و فمینیسم که استدلال می‌کند:
\begin{itemize}[nosep]
    \item پدرسالاری قدیمی‌ترین شکل سلطه است
    \item دولت و سرمایه‌داری پدرسالارانه‌اند
    \item آنارشیسم بدون فمینیسم ناقص است
    \item فمینیسم بدون نقد دولت و سرمایه ناکافی است
\end{itemize}
\textbf{پیشگام:} \person{اما گلدمن} (۱۸۶۹-۱۹۴۰)
\end{defbox}

\begin{personbox}[اما گلدمن (۱۸۶۹-۱۹۴۰)]

\textbf{ملیت:} روسی-آمریکایی

\textbf{لقب:} «زن خطرناک‌ترین آمریکا»

\textbf{آثار:} \book{آنارشیسم و مقالات دیگر}، \book{زیستن زندگی من}

\mediumspace

\person{گلدمن} آنارشیسم را به مسائل زنان، جنسیت و زندگی روزمره پیوند زد:

\begin{itemize}[nosep]
    \item دفاع از کنترل زاد و ولد (غیرقانونی بود)
    \item نقد ازدواج سنتی به عنوان «بردگی»
    \item دفاع از آزادی جنسی
    \item «اگر نتوانم برقصم، این انقلابِ من نیست»
    \item نقد انقلاب روسیه از نزدیک
\end{itemize}
\end{personbox}

\subsection{آنارشیسم سبز و بدوی‌گرایی}

\begin{defbox}[تعریف: اکو-آنارشیسم]
پیوند آنارشیسم با اکولوژی:
\begin{itemize}[nosep]
    \item بحران محیط زیست ناشی از سلطه بر طبیعت است
    \item سلطه بر طبیعت و سلطه بر انسان به هم پیوسته‌اند
    \item جامعه آنارشیست باید پایدار و اکولوژیک باشد
\end{itemize}
\textbf{چهره کلیدی:} \person{ماری بوکچین} (اکولوژی اجتماعی)
\end{defbox}

\begin{defbox}[تعریف: آنارکو-بدوی‌گرایی]
گرایش رادیکال‌تر که استدلال می‌کند:
\begin{itemize}[nosep]
    \item تمدن خود مشکل است، نه فقط دولت
    \item کشاورزی، صنعت، شهر—همه سرکوبگرند
    \item باید به زندگی شکارچی-گردآورنده بازگشت
\end{itemize}
\textbf{چهره کلیدی:} \person{جان زرزان}

\textbf{نقد:} آیا این عملی است؟ آیا میلیاردها انسان می‌توانند شکارچی-گردآورنده باشند؟
\end{defbox}

\subsection{پست-آنارشیسم}

\begin{defbox}[تعریف: پست-آنارشیسم]
ترکیب آنارشیسم با نظریه پساساختارگرا (فوکو، دریدا، دلوز):
\begin{itemize}[nosep]
    \item تمرکز بر \textbf{قدرت پراکنده}، نه فقط دولت
    \item نقد «ذات‌گرایی» در آنارشیسم کلاسیک
    \item توجه به زبان، گفتمان و سوبژکتیویته
    \item آنارشیسم به عنوان \textbf{اتیک}، نه برنامه سیاسی
\end{itemize}
\textbf{چهره‌ها:} سیمون نیومن، لوئیس کال
\end{defbox}

% ══════════════════════════════════════════════════════════════════════════════
\section{نقد آنارشیسم}
% ══════════════════════════════════════════════════════════════════════════════

\subsection{نقد از چپ مارکسیستی}

\begin{quotebox}[فردریش انگلس، درباره اقتدار (۱۸۷۲)]
\textit{«آیا این آقایان هرگز انقلابی دیده‌اند؟ انقلاب مسلماً اقتدارگراترین چیزی است که وجود دارد.»}
\end{quotebox}

\begin{itemize}[nosep, rightmargin=1.5cm]
    \item آنارشیسم «آرمان‌شهرگرا» و غیرعملی است
    \item بدون دولت کارگری، نمی‌توان در برابر ضدانقلاب مقاومت کرد
    \item تاریخ نشان داده: آنارشیست‌ها همیشه شکست خورده‌اند
    \item سازمان‌دهی افقی در مقیاس بزرگ ممکن نیست
\end{itemize}

\subsection{نقد از راست}

\begin{itemize}[nosep, rightmargin=1.5cm]
    \item انسان ذاتاً نیازمند نظم و سلسله‌مراتب است
    \item بدون دولت، زورمندان حاکم می‌شوند
    \item مالکیت خصوصی حق طبیعی است
    \item آنارشیسم به هرج‌ومرج و خشونت می‌انجامد
\end{itemize}

\subsection{پاسخ آنارشیست‌ها}

\begin{itemize}[nosep, rightmargin=1.5cm]
    \item «دولت کارگری» شوروی به چه انجامید؟ پیش‌بینی باکونین درست درآمد
    \item انقلاب اسپانیا نشان داد آنارشیسم عملی است—سرکوب خارجی شکستش داد
    \item سازمان‌دهی افقی امروز با فناوری ممکن‌تر است
    \item بسیاری از جوامع بدون دولت در تاریخ وجود داشته‌اند
\end{itemize}

% ══════════════════════════════════════════════════════════════════════════════
\section{آنارشیسم امروز}
% ══════════════════════════════════════════════════════════════════════════════

\begin{personbox}[نوام چامسکی (متولد ۱۹۲۸)]

\textbf{ملیت:} آمریکایی

\textbf{حوزه‌ها:} زبان‌شناسی، فلسفه، فعالیت سیاسی

\textbf{گرایش:} آنارکو-سندیکالیسم

\mediumspace

\person{چامسکی}، زبان‌شناس برجسته و فعال سیاسی، خود را آنارشیست می‌داند:

\begin{itemize}[nosep]
    \item هر شکل از اقتدار باید خود را توجیه کند
    \item اگر نتوانست، باید برچیده شود
    \item نقد سیاست خارجی آمریکا
    \item نقد رسانه‌های اصلی («ساخت رضایت»)
    \item اما: «آنارشیسم من آرمان دوردست است؛ امروز باید از دستاوردها دفاع کرد»
\end{itemize}
\end{personbox}

\begin{historybox}[آنارشیسم در جنبش‌های معاصر]
آنارشیسم در جنبش‌های اخیر حضور پررنگی داشته:
\begin{itemize}[nosep]
    \item \textbf{آلترگلوبالیسم} (سیاتل ۱۹۹۹): تاکتیک‌های آنارشیستی
    \item \textbf{اشغال وال‌استریت} (۲۰۱۱): سازمان‌دهی افقی، بدون رهبر
    \item \textbf{روژاوا} (کردستان سوریه): الهام از بوکچین، «کنفدرالیسم دموکراتیک»
    \item \textbf{زاپاتیست‌ها} (مکزیک): خودگردانی بومی
    \item \textbf{بلوک سیاه}: تاکتیک‌های تظاهراتی
\end{itemize}
\end{historybox}

% ══════════════════════════════════════════════════════════════════════════════
%                         خلاصه و پرسش‌ها
% ══════════════════════════════════════════════════════════════════════════════

\newpage

\begin{summarybox}

\textbf{۱. تعریف:} آنارشیسم نه هرج‌ومرج، بلکه نفی سلطه و سلسله‌مراتب اجباری است. هدف: خودگردانی و همکاری داوطلبانه.

\textbf{۲. پیشگامان:} گادوین (عقل‌گرایی)، پرودون («مالکیت دزدی است»)، اشتیرنر (فردگرایی افراطی).

\textbf{۳. آنارشیسم جمعی:} باکونین—انقلاب همزمان علیه دولت و سرمایه. جدال با مارکس.

\textbf{۴. آنارکو-کمونیسم:} کروپوتکین—کمک متقابل، توزیع بر اساس نیاز، علم و آنارشیسم.

\textbf{۵. آنارکو-سندیکالیسم:} اتحادیه‌ها ابزار مبارزه و اداره جامعه. انقلاب اسپانیا بزرگ‌ترین تجربه.

\textbf{۶. آنارشیسم فردگرا:} تاکر و سنت آمریکایی. مناقشه با آنارکو-کاپیتالیسم.

\textbf{۷. جریان‌های معاصر:} آنارکا-فمینیسم، اکو-آنارشیسم، پست-آنارشیسم.

\textbf{۸. امروز:} آنارشیسم در جنبش‌های ضد جهانی‌سازی، اشغال، روژاوا حضور دارد.

\end{summarybox}

\vspace{20pt}

\begin{questionbox}

\begin{enumerate}[nosep, rightmargin=1cm]
    \item آیا پیش‌بینی باکونین درباره دولت کارگری (که به بوروکراسی تبدیل می‌شود) درست درآمد؟

    \item انقلاب اسپانیا چرا شکست خورد؟ آیا این شکست ناشی از ضعف ذاتی آنارشیسم بود یا عوامل بیرونی؟

    \item آیا آنارکو-کاپیتالیسم واقعاً «آنارشیسم» است؟ چرا آنارشیست‌های سنتی آن را رد می‌کنند؟

    \item با فناوری‌های امروز (اینترنت، رمزارز، سازمان‌دهی شبکه‌ای)، آیا آنارشیسم عملی‌تر شده است؟
\end{enumerate}

\end{questionbox}

% ══════════════════════════════════════════════════════════════════════════════
%                         پایان فصل ششم
% ══════════════════════════════════════════════════════════════════════════════